\documentclass[11pt,a4paper,oneside,fontset=none,scheme=plain]{ctexbook}

\usepackage{lmodern}
\setCJKmainfont{Noto Serif CJK SC}[BoldFont={Noto Serif CJK SC},BoldFeatures={FakeBold=2},ItalicFont={Noto Serif CJK SC},ItalicFeatures={FakeSlant=.15}]
\setCJKsansfont{Noto Sans CJK SC}
\setCJKmonofont{Noto Sans Mono CJK SC}
\renewcommand{\contentsname}{目录}
\renewcommand{\figurename}{图}
\renewcommand{\tablename}{表}
\renewcommand{\bibname}{参考文献}
% 每篇独立编号；无编号的篇题不应让section/subsection出现零章前缀。
\renewcommand{\thesection}{\arabic{section}}
\renewcommand{\proofname}{证明}
\ctexset{chapter={format=\huge\bfseries,beforeskip=50pt,afterskip=40pt},section={format=\Large\bfseries},subsection={format=\large\bfseries}}
\usepackage[a4paper,inner=30mm,outer=27mm,top=28mm,bottom=30mm]{geometry}
\usepackage{amsmath,amssymb,amsthm,mathrsfs,mathtools,stmaryrd}
\usepackage{tikz,tikz-cd}
\usetikzlibrary{
  arrows,
  arrows.meta,
  calc,
  decorations.pathmorphing,
  matrix,
  positioning,
  shapes,
  shapes.geometric
}
\usepackage{graphicx,float,adjustbox,enumitem,multicol}
\usepackage{etoolbox}
% 标准 samepage 只定义入口；给环境结束钩子提供空出口，兼容 ctex/LaTeX 环境钩子。
\ifdefined\endsamepage\else\def\endsamepage{}\fi
\usepackage{xcolor}
\usepackage[hidelinks,hypertexnames=false,pdfstartview={},linktoc=page,
  pdftitle={SGA 3 群概形：简体中文原生重建工作稿},
  pdfauthor={M. Demazure; A. Grothendieck et al.},
  pdfsubject={SGA 3 第一部分；覆盖范围详见 PROGRESS.json},
  pdfkeywords={SGA 3, 群概形, 中文, 原生 LaTeX}]{hyperref}

\setlength{\emergencystretch}{3em}
\raggedbottom

\theoremstyle{plain}
\newtheorem{theorem}{定理}[section]
\newtheorem{proposition}[theorem]{命题}
\newtheorem{corollary}[theorem]{推论}
\newtheorem{lemma}[theorem]{引理}
\newtheorem{scholium}[theorem]{评注}
\newtheorem{application}[theorem]{应用}
\newtheorem{complements}[theorem]{补充}

\theoremstyle{definition}
\newtheorem{definition}[theorem]{定义}
\newtheorem{example}[theorem]{例}

\theoremstyle{remark}
\newtheorem{remark}[theorem]{注}

\providecommand{\Ga}{\mathbb G_a}
\providecommand{\Gm}{\mathbb G_m}
\providecommand{\Norm}{\operatorname{Norm}}
\providecommand{\Centr}{\operatorname{Centr}}
\providecommand{\Ker}{\operatorname{Ker}}
\providecommand{\Dom}{\operatorname{Dom}}
\providecommand{\rank}{\operatorname{rank}}
\providecommand{\Spec}{\operatorname{Spec}}
\providecommand{\Lie}{\operatorname{Lie}}
\providecommand{\intaut}{\operatorname{int}}
\providecommand{\Fib}{\operatorname{Fib}}
\providecommand{\Aut}{\operatorname{Aut}}
\providecommand{\Autext}{\operatorname{Autext}}
\providecommand{\rgred}{\operatorname{rgred}}
\providecommand{\rgss}{\operatorname{rgss}}
\providecommand{\GL}{\operatorname{GL}}
\providecommand{\Qp}{\mathbb{Q}_{+}}
\providecommand{\cR}{\mathcal{R}}
\providecommand{\VR}{\mathcal{V}}
\providecommand{\Nil}{\operatorname{Nil}}

\providecommand{\SGAParagraph}[1]{%
  \par\medskip\noindent\textbf{#1.}---\quad
}
\providecommand{\NDE}[1]{%
  \footnote{\textit{编者注.} #1}%
}
\newenvironment{sgastatement}[1]
  {\par\medskip\noindent\textbf{#1.}---\ \itshape}
  {\par\medskip}
\providecommand{\SGAQED}{\hfill\(\square\)\par\smallskip}
\providecommand{\Idx}[2]{%
  \par\noindent\hangindent=1.2em\hangafter=1
  #1\if\relax\detokenize{#2}\relax\else,\ #2\fi\par
}
\providecommand{\SGASection}[2]{%
  \section*{#1.\ #2}%
  \addcontentsline{toc}{section}{#1.\ #2}%
}
\providecommand{\SGASubsection}[2]{%
  \subsection*{#1.\ #2}%
  \addcontentsline{toc}{subsection}{#1.\ #2}%
}
\providecommand{\SGAStatement}[2]{%
  \par\medskip\noindent\textbf{#1.}---\ {\itshape #2\par}\medskip
}
\providecommand{\SGAProof}[1]{%
  \par\smallskip\noindent\textit{#1.}\quad
}

\makeatletter
\providecommand{\SGARefTarget}[1]{\phantomsection\label{#1}}
\providecommand{\SGAInlineRefTarget}[1]{%
  \begingroup
    \Hy@raisedlink{\hyper@anchorstart{#1}\hyper@anchorend}%
    \def\@currentHref{#1}%
    \label{#1}%
  \endgroup}
\providecommand{\SGAInlineEquationTarget}[1]{%
  \begingroup
    \Hy@raisedlink{\hyper@anchorstart{#1}\hyper@anchorend}%
    \def\@currentHref{#1}%
    \ltx@label{#1}%
  \endgroup}
\let\SGA@PendingSectionTarget\@empty
\providecommand{\SGANextSectionTarget}[1]{%
  \gdef\SGA@PendingSectionTarget{#1}}
\providecommand{\SGA@UsePendingSectionTarget}{%
  \ifx\SGA@PendingSectionTarget\@empty\else
    \expandafter\SGAInlineRefTarget\expandafter{%
      \SGA@PendingSectionTarget}%
    \global\let\SGA@PendingSectionTarget\@empty
  \fi}
\renewcommand{\SGASection}[2]{%
  \section*{\SGA@UsePendingSectionTarget#1.\ #2}%
  \addcontentsline{toc}{section}{#1.\ #2}%
}
\renewcommand{\SGASubsection}[2]{%
  \subsection*{\SGA@UsePendingSectionTarget#1.\ #2}%
  \addcontentsline{toc}{subsection}{#1.\ #2}%
}
\pdfstringdefDisableCommands{%
  \def\SGARefLink#1#2{#2}%
  \def\SGARefPending#1#2{#2〔引文待核〕}%
  \def\SGARefExternal#1#2{#2}%
  \def\SGARefTarget#1{}%
  \def\SGAInlineRefTarget#1{}%
  \def\SGAInlineEquationTarget#1{}%
  \def\SGANextSectionTarget#1{}%
}
\providecommand{\SGARefLink}[2]{%
  \@ifundefined{r@#1}{#2}{\hyperref[#1]{#2}}%
}
% 有界引用校订：疑点保留原显示文字并明标待核；外引不跳入本书同号条目。
\DeclareRobustCommand{\SGARefPending}[2]{%
  #2\ifmmode\text{\normalfont\scriptsize〔引文待核〕}%
  \else{\normalfont\scriptsize〔引文待核〕}\fi}
\DeclareRobustCommand{\SGARefExternal}[2]{#2}
\providecommand{\GuideChapter}[2]{%
  \@dottedtocline{0}{0em}{0em}{#1}{#2}}
\providecommand{\GuideChapterSplit}[3]{%
  \par\noindent\textbf{#1}\par\nobreak
  \@dottedtocline{0}{2em}{0em}{#2}{#3}}
\providecommand{\GuideSection}[2]{%
  \@dottedtocline{1}{1.45em}{2.25em}{#1}{#2}}
\newcommand{\SGAInputRoot}[1]{%
  \def\input@path{{inputs/#1/}}%
}
\newcommand{\SGATagsLeft}{%
  \tagsleft@true
  \let\veqno\@@leqno
}
\newcommand{\SGATagsRight}{%
  \tagsleft@false
  \let\veqno\@@eqno
}
\makeatother

\newcommand{\SGAStartExpose}[1]{%
  \clearpage
  \setcounter{section}{0}%
  \setcounter{subsection}{0}%
  \setcounter{theorem}{0}%
  \setcounter{footnote}{0}%
  \setcounter{figure}{0}%
  \SGAInputRoot{#1}%
}

